\section{Rationale}

Relational control is a working label for a candidate person-level propensity to intentionally initiate changes in contact, availability, boundaries, or interaction structure. Its stability and distinctness are empirical questions. The operational definition does not require a discrepancy equation, successful regulation, or a particular controller family. Relational goals may motivate an action without establishing the ARC discrepancy-to-action hypothesis.

The construct is not equated with dominance, aggression, coercion, attachment anxiety, avoidance, general agency, assertiveness, or a dynamical gain parameter. Consent, perceived effectiveness, and harmfulness of an action are separate properties, not interchangeable indicators of this candidate propensity. Frequent attempts can reflect more opportunities or stressors rather than a stronger person-level tendency.

This manuscript owns the construct and measurement question. The separate identifiability manuscript owns discrimination among synthetic dynamical generators. Recovering a generator cannot validate this questionnaire; a reliable questionnaire cannot identify a psychological P, PI, or PID mechanism. An empirical relation between scale scores and dynamic parameters would require a separately planned bridge study.
